\chapter{泛函积分}
\renewcommand{\currentchapterpath}{chapters/ch06/}

\begin{flushright}
\textit{“形而上者谓之道，形而下者谓之器。” ——《周易·系辞》
\footnote{路径积分把算符“器”化为轨迹上的作用量“道”，是多体场论语言的入口。}}
\end{flushright}

泛函积分（functional integral）是相干态表象下的路径积分：把迹与时间演化写成对玻色/Grassmann 场组态的积分。本节按施均仁讲义第 3 章给出骨架，并与 Matsubara 格林函数衔接。


\chapterinput{feynman_path}
\chapterinput{imaginary}
\chapterinput{coherent_path}
\chapterinput{partition_gf}
